% !TeX root = ../despliegue.tex
% ===========================================================================
%  Comprobaciones tras el arranque.
% ===========================================================================

\section{Verificar que levantó bien}

Tres comprobaciones bastan:

\begin{comando}
docker compose ps                                       # los once en healthy
curl -u cliente1:cliente1 http://localhost/api/canchas   # el catálogo
open http://localhost                                    # la aplicación
\end{comando}

La segunda es la más informativa de las tres: si devuelve el catálogo, entonces
el edge, el gateway, \servicio{ms-usuarios} —que verificó la credencial—,
\servicio{ms-canchas} y PostgreSQL están todos funcionando y comunicándose entre
sí.

Una petición sin credenciales debe responder \cmd{401}:

\begin{comando}
curl -i http://localhost/api/canchas
\end{comando}

El Swagger de cada microservicio queda accesible por su puerto directo:

\begin{itemize}
  \item \url{http://localhost:8081/swagger-ui.html} — usuarios y autenticación
  \item \url{http://localhost:8082/swagger-ui.html} — canchas
  \item \url{http://localhost:8083/swagger-ui.html} — reservas
  \item \url{http://localhost:8084/swagger-ui.html} — reportes
\end{itemize}
